\section{DeepSeek, código abierto y código cerrado}

La sección anterior sobre la Scaling Law trató el crecimiento de las capacidades; esta sección pasa a cómo la industria percibe esas capacidades y cómo se difunden. DeepSeek llevó el código abierto, la eficiencia y los resultados reales al centro del discurso de la industria.

En 2025, el título de esta sección fue «Aprender de DeepSeek». Zhang Peng (张鹏) mencionó que, en cierta medida, DeepSeek también le enseñó algo a la industria: en lo técnico no siempre hacen falta grandes eventos de lanzamiento, porque al final lo que cuenta es el resultado en aplicaciones reales. DeepSeek cambió a la vez el discurso sobre el código abierto y el discurso sobre la eficiencia, y llevó a las empresas de modelos a replantear la relación entre apertura, ecosistema y control comercial.

\begin{figure}[H]
\centering
\includegraphics[width=0.94\textwidth]{open-vs-closed.png}
\caption{Código abierto frente a código cerrado: el equilibrio entre la influencia en el ecosistema y el control comercial.}
\end{figure}

\begin{knowledgebox}{Lectura de la figura: el código abierto y el código cerrado no son una dicotomía moral}
El código abierto permite ampliar el ecosistema, recibir retroalimentación de los desarrolladores y construir una reputación técnica; el código cerrado permite conservar el control comercial, atender a clientes empresariales y manejar los límites de seguridad y cumplimiento normativo. Por lo general, las empresas de modelos necesitan combinar ambos enfoques en distintas capas de producto y distintos segmentos de clientes.
\end{knowledgebox}

\subsection{La retroalimentación global que trae el código abierto}

Zhang Peng mencionó que, después de que Zhipu (智谱) liberó sus modelos como código abierto, la evaluación en el extranjero fue muy positiva, e incluso hubo otros productos o empresas que usaron, destilaron o recortaron directamente los modelos. Esto muestra que el código abierto tiene influencia, pero también que da lugar a un ecosistema complejo en el que los modelos se reutilizan, se modifican y se vuelven a empaquetar.

\begin{warningbox}{El código abierto no solo trae elogios}
El código abierto aporta reputación entre los desarrolladores y difusión en el ecosistema, pero también trae problemas como la atribución comercial difusa, el reempaquetado de los modelos por terceros y la dificultad para capturar ingresos. La estrategia de código abierto debe diseñarse junto con la estrategia de comercialización.
\end{warningbox}

\subsection{Resumen de la sección}

Después de DeepSeek, la competencia entre empresas de modelos da más peso a la eficiencia, a los resultados reales y al ecosistema de código abierto. El código abierto y el código cerrado no son opciones excluyentes: la cuestión es cómo cada empresa de modelos reparte la influencia en el ecosistema, el valor comercial y los límites de seguridad.
